% 第三问摘要段落，供并入总论文的 abstract 环境。
% 本文件为正文片段，无需另加摘要标题、abstract 环境或关键词。
% 以下起报提前量均为名义值；预警提前时间以核实启动时刻计算。

\textbf{对于问题三，}建立双机制风险外推、集合情景估计与分段传播相结合的预测模型。沿用第二问的热力—冰冻圈和降雨—水分指标，引入GEFS集合均值与离散度，通过天气和融水参数扰动构造1000个合成情景，以持续越阈识别预警起点，以持续风险峰值及有效情景的分位数估计触发时刻和区间。同一事件的五次滚动回放表明，原始95\%情景区间随提前量缩短由155 h收窄至39 h，覆盖核实启动时刻4次，描述性覆盖比例为80\%。24 h与48 h起报的触发中位时刻误差分别为0.13 h和2.87 h，橙色预警分别提前21.87 h和14.87 h。进一步结合灾后反演速度建立A—B—C—D分段传播模型，24 h起报下源区至口岸的传播时间中位数为7.39 min，95\%情景范围为7.12—7.70 min，既定沟道的地形修正长度为21.681 km。传播时间的波动为分钟量级，口岸到达区间主要取决于源区触发时间的不确定性。
\par
